\documentclass{article}

\usepackage[utf8]{inputenc}
\usepackage[T1]{fontenc}
\usepackage[english]{babel}
\usepackage{microtype}

\usepackage{amsmath, amssymb, amsthm}
\usepackage{mathtools}
\usepackage{totcount}
\usepackage{enumitem}

\usepackage{newtxtext,newtxmath}

\usepackage{titlesec}
\titleformat{\chapter}[display]
  {\normalfont\huge\bfseries} % Format of the text
  {Lecture \thechapter}       % Label (e.g., "Lecture 1")
  {0.5em}                     % Spacing between label and title
  {\Huge}                     % Format of the title

% Theorem-like environments
\newtheorem{theorem}{Theorem}[section]
\newtheorem{lemma}[theorem]{Lemma}
\newtheorem{proposition}[theorem]{Proposition}
\newtheorem{corollary}[theorem]{Corollary}

\theoremstyle{definition}
\newtheorem{definition}[theorem]{Definition}
\newtheorem{example}[theorem]{Example}

\theoremstyle{remark}
\newtheorem{remark}[theorem]{Remark}

\title{Practice Midterm 2 for Elementary Number Theory}
\author{Egbert Rijke}
\date{\today}

\usepackage{hyperref}
\hypersetup{
    colorlinks=true,
    linkcolor=purple,
    filecolor=magenta,
    urlcolor=blue,
    citecolor=purple,
    pdftitle={Lecture notes on Number Theory},
    pdfpagemode=FullScreen,
}
\usepackage[capitalize,english]{cleveref}

  % We set up the exercise environment, which produces list environment in a new
% unnumbered subsection that also gets mentioned in the table of contents.

\newtotcounter{exercisecounter}
\newcommand{\exercise}{\item\stepcounter{exercisecounter}}

\newlist{exenum}{enumerate}{1}
\setlist[exenum]{noitemsep,label=\arabic*}
  %,ref=\thechapter.\arabic*}

\crefname{exenumi}{Exercise}{Exercises}
  
\newlist{subexenum}{enumerate}{1}
\setlist[subexenum]{noitemsep,label=(\alph*),ref=\theexenumi.\alph*}
\crefname{subexenumi}{Exercise}{Exercises}
  
\newenvironment{exercises}
{%
\subsection*{Exercises}%
\addcontentsline{toc}{section}{Exercises}%
\sectionmark{Exercises}%
\begin{exenum}}
{%
%\addtocounter{exercisecounter}{\theexenumi}%
\end{exenum}}

\begin{document}
\maketitle

\begin{exercises}
  \exercise Solve the following linear congruences:
  \begin{subexenum}
  \item $3x\equiv 7\pmod{11}$,
  \item $12x\equiv 18\pmod{30}$,
  \item $7x\equiv 14\pmod{21}$.
  \end{subexenum}
  Describe the congruence classes of all solutions, if multiple solutions are possible.
  \exercise Find the smallest positive integer $x$ such that
  \begin{align*}
    x & \equiv 2 \pmod{3} \\
    x & \equiv 5 \pmod{7} \\
    x & \equiv 6 \pmod{11}.
  \end{align*}
  \exercise Determine the residue class of
  \begin{equation*}
    2^{100}+3^{100}
  \end{equation*}
  modulo $13$.
  \exercise Determine the residue class of
  \begin{equation*}
    7^{4^{50}}
  \end{equation*}
  modulo $13$.
  \exercise
  \begin{subexenum}
  \item Use Legendre's formula or Kummers theorem (your choice) to compute the prime factorization of
    \begin{equation*}
      \binom{18}{6}.
    \end{equation*}
  \item Determine $\sigma\binom{18}{6}$ and $\phi\binom{18}{6}$, where $\sigma$ is the sum-of-divisors function and $\phi$ is Euler's totient function.
  \end{subexenum}
\end{exercises}

\end{document}
